%!TEX root = ../thesis-main.tex
\section{Dependency Management in Simulation Algorithms}

The discrete-event paradigm reduces simulation to repeatedly selecting the next
scheduled event, applying a state transition and updating the set of future
scheduled events. In practice, the computational cost of this loop is often
dominated not by the event handlers themselves, but by their \emph{dependency
	management}, namely \emph{what} must be recomputed after a state change and
\emph{which} future events are affected.

This section focuses on algorithmic mechanisms that exploit sparsity in those
dependencies to reduce overhead. We first discuss the \ac{SSA}, where
dependencies arise from reaction stoichiometry and only a small subset of
propensities must be updated after each reaction. Then we broaden the notion of
dependency to \emph{causal dependencies} introduced by parallel and distributed
execution: once the event list and the state are partitioned across multiple
workers, simulation correctness depends on enforcing a strict causal ordering
and execution of events.  [CONTINUES...]

\subsection{Manual Dependency Tracking: The Stochastic Simulation Algorithm}\label{ssec:mdt}
While general \ac{DES} frameworks often rely on empirical distributions for
event timings, the \ac{SSA}~\cite{GillespieESSCCR:1977} provides a physically
grounded method for simulating the time-evolution of chemical reaction networks.
The "Direct Method" of the \ac{SSA} relies on the \emph{Markovian} property:
the probability of a reaction occurring depends solely on the current
molecular population (the state) rather than the history of the system.

In this framework, each reaction $R_\mu$ is assigned a \emph{propensity function}
$a_\mu$. Assuming the system is well-stirred and at constant volume, the
time to the next reaction $\tau$ is an exponential random variable with
a cumulative rate equal to the aggregate propensity $\lambda = \sum_j a_j$.
The \ac{PDF} of this timing is given by:

\begin{equation}\label{eq:nextreactiontime}
	P(\tau) = \lambda e^{-\lambda \tau}
\end{equation}

The specific reaction $\mu$ that occurs at that time is a discrete random variable
chosen with probability:

\begin{equation}\label{eq:nextreactionprobability}
	P(\text{reaction} = \mu) = \frac{a_{\mu}}{\lambda}
\end{equation}

To simulate this, two independent random numbers $r_1, r_2 \in (0, 1]$ are
sampled from a uniform distribution. The time step is then computed via the
\emph{Inverse Transform Sampling} as $\tau = -\ln(r_1)/\lambda$, while $\mu$ is
the smallest integer satisfying $\sum_{j=1}^{\mu} a_j > r_2 \lambda$.

\subsubsection{\ac{SSA} Optimisation}

The \ac{SSA} can be then summarised in \Cref{alg:ssa}. Right away some performance
issues arise by looking at the pseudocode:
\begin{itemize}
	\item Every propensity $a_j$ is recomputed even if the reaction did not
	      involve those species;
	\item All propensities have to be summed all together at every step;
	\item If reactions are stored in a list-like structure, finding the
	      next reaction $\mu$ to execute is done with linear complexity.
\end{itemize}

\begin{algorithm}
	\caption{Stochastic Simulation Algorithm (Direct Method)}\label{alg:ssa}
	\begin{algorithmic}[1]
		\Procedure{SSA}{initial state $\mathbf{x}$, time $T_{end}$}
		\State $t \gets 0$
		\While{$t < T_{end}$}
		\For{each reaction $R_j$}
		\State Calculate propensity $a_j$ based on $\mathbf{x}$
		\EndFor
		\State $\lambda \gets \sum a_j$
		\State Sample $r_1, r_2 \sim \text{Uniform}(0, 1)$
		\State $\tau \gets \frac{-\ln(r_1)}{\lambda}$ \Comment{Time to next event}
		\State Find smallest $\mu$ such that $\sum_{j=1}^{\mu} a_j > r_2 \lambda$ \Comment{Select event}
		\State Update state $\mathbf{x}$ according to reaction $R_\mu$
		\State $t \gets t + \tau$
		\EndWhile
		\EndProcedure
	\end{algorithmic}
\end{algorithm}

In short, \ac{SSA} complexity is $O(M)$ with $M$ being the number of reactions
in the system. In addition to that, the direct method takes time proportional
to the number of reactions to calculate $\lambda$ and to generate a random number
according to \Cref{eq:nextreactionprobability}.

\paragraph{Optimising propensity updates}
It is possibile to optimise the execution of such algorithm by noting that,
in most systems, chemical reactions are \emph{sparse}, hence only a small
subset of propensities may change from step to step. If a reaction
$A$ consumes species $X$, only those concentration that also use $X$ will be
affected. It is possible to model such influence relationships among reactions
by means of a \emph{Dependency Graph} $\mathcal{G} = (V, E)$ where:
\begin{itemize}
	\item $V$ is the set of reactions;
	\item An edge $e_{i,j} \in E$ exists from $R_i \in V$ to $R_j \in V$ if the
	      execution of the reaction $R_i$ changes the set of substances involved
	      in $R_j$.
\end{itemize}
With this data structure we cut the cost of updating $M = |V|$ propensities at
every iteration with just a small constant number of them.

\paragraph{Optimising scheduling}
As already said, a major bottleneck in simulating such systems is the \ac{FEL}
management: in both the direct method and in the \emph{First Reaction
	Method}~\cite{GILLESPIE1976403} linear scans are performed to find the
next reaction $\mu$ to execute, where in the latter the reaction with
the smallest \emph{putative time} $\tau_j$ should be chosen.
Starting from this idea, in ~\cite{GibsonBruckEESSCSMSMC:2000} the \emph{Next Reaction
	Method} is presented by introducing four important properties:
\begin{enumerate}
	\item The propensity function $a_j$ is stored along with the putative time
	      $\tau_j$. If $a_j$ is not recomputed, nor does $\tau_j$;
	\item A dependency graph $\mathcal{G} = (V, E)$ is stored which will help
	      telling precisely which set of reactants to change when a given
	      reaction is executed.
	\item The list-like data structure backing the \ac{FEL} is replaced with an
	      \ac{IPQ}: it is a binary heap combined with a map,
	      which maximises efficiency for sparse operations and a small number of
	      updates, granting $O(1)$ access and a sorting cost of $O(\log M)$.
	\item \emph{Random number reuse}: This property addresses the computational
	      cost of generating high-quality random numbers (could be the most heavy
	      task in simulation). In the direct method, every step requires two new
	      random numbers. In the Next Reaction Method, if a reaction $R_j$ is
	      \emph{not} affected by the execution of $R_\mu$, its absolute putative
	      time $T_j$ is simply kept in the \ac{FEL}. However, if $R_j$ is
	      affected (its propensity changes from $a_j$ to $a_j^{new}$), its new
	      putative time $T_j^{new}$ is rescaled to maintain statistical
	      exactness:
	      \begin{equation}\label{eq:rand_reuse}
		      T_j^{new} = t + \frac{a_j}{a_j^{new}} (T_j^{old} - t)
	      \end{equation}
	      This transformation preserves the Markovian property because the
	      exponential distribution is memoryless. Therefore, a new random number
	      is only required for the reaction that just fired, while all others are
	      updated algebraically.
\end{enumerate}
With the above elements, the \emph{Next Reaction Method} transform a linear,
brute-force search into a targeted, logarithmic update cycle. Through the
Dependency Graph it is possibile to identify which reactions in the \ac{IPQ}
need to be modified when a reaction executes. With the latter it is possibile
to perform a lookup of a given reaction in constant time, and update its
putative time using \Cref{eq:rand_reuse} which drastically reduce the time
spent in computing new random values. Finally, the heap structure of the
\ac{IPQ} restores the chronological order of the \ac{FEL} in logarithmic time,
ensuring the positioning of the next event at the root of the tree.

\begin{algorithm}
	\caption{Next Reaction Method (Gibson-Bruck)}\label{alg:gibson_bruck}
	\begin{algorithmic}[1]
		\Procedure{NextReactionMethod}{initial state $\mathbf{x}$, dependency graph $\mathcal{G}$}
		\State $t \gets 0$
		\State Generate $a_j$ and putative times $T_j$ for all $R_j$ \Comment{Initial expensive step}
		\State Build \texttt{\ac{IPQ}} containing all $T_j$
		\While{$t < T_{end}$}
		\State $(\mu, T_{\mu}) \gets \texttt{\ac{IPQ}.top()}$ \Comment{$O(1)$ access to next event}
		\State $t \gets T_{\mu}$ \Comment{Advance CLOCK to absolute event time}
		\State Update state $\mathbf{x}$ according to reaction $R_{\mu}$
		\State Sample new $r$ and update $T_{\mu} = t - \frac{\ln(r)}{a_{\mu}}$ \Comment{New time only for reaction that fired}

		\For{each reaction $R_j$ in $\mathcal{G}$ affected by $R_{\mu}$}
		\If{$j \neq \mu$}
		\State $a_j^{new} \gets \text{update propensity using } \mathbf{x}$
		\State $T_j \gets t + \frac{a_j^{old}}{a_j^{new}}(T_j - t)$ \Comment{Reuse/rescale random numbers}
		\State \texttt{\ac{IPQ}.update}($j, T_j$) \Comment{$O(\log M)$ heap maintenance}
		\EndIf
		\EndFor
		\EndWhile
		\EndProcedure
	\end{algorithmic}
\end{algorithm}

A simplified version of the algorithm is then presented in
\Cref{alg:gibson_bruck}, which shows that for each step we pay $O(\log M)$ for
the heap updates and $O(\text{deg}(\mathcal{G}))$ for the dependency graph
traversal.

\subsection{Synchronisation-based Approaches}

The event scheduling algorithm discussed so far assumes a \emph{centralised}
execution mechanism, where a single simulation engine maintains the global
\ac{FEL} and processes events in non-decreasing order. It is easy to assess
this design's correctness, however it quickly becomes a scalability bottleneck
as the number of model entities and scheduled events grows.
\ac{PDES} addresses this problem by partitioning the model into a set of
interacting components, each executed by a separate worker (a thread, a
physical core or a networked process). With the common abstraction of \ac{LP}
we indicate a computing owning a portion of the model state, maintaining its
local \ac{FEL} while interacting with other \ac{LP}s by exchanging messages.
The simulation output of a parallel simulator should be the same as executing
that model under the same assumptions with a sequential simulator, with the
only difference being that parallel and distributed simulation should allow
multiple LPs to run concurrently in order to gain some sort of computational
speedup by the employment of multiple resources in
parallel~\cite{fujimoto1993FAPDESWFS}.

\paragraph{Causality and Simulation Correctness}
In general, when one component may directly or indirectly influence another
component, there is a causal dependency of events. To guarantee execution
correctness, the generation of
states must preserve what is called \emph{causality}. Moreover,
~\cite{Fujimoto1990PDES} states that correct simulation happens under the
\emph{local causality constraint}, namely when each \ac{LP} processes the
events in timestamp order.
In order to avoid \emph{synchronisation problems}, every \ac{LP} should adhere
to the local causality constraint.
The output of a parallel simulator is equivalent to its sequential counterpart
if and only if the local causality constraint is respected. Despite this
observation, in presence of independent events their out-of-order execution
will not cause causality errors. This raises an issue linked to the way
parallel simulators can determine whether two events could be executed
concurrently, and whether or not \ac{LP}s are allowed to temporarily violate
the local causality constraint. Two opposite approaches has been developed
during the years and will be presented in the following sections.

\subsubsection{Conservative Execution}\label{sssec:conservative_pdes}

\emph{Conservative execution} is a form of processes synchronisation which
maintains correctness of a simulation by ensuring that an \ac{LP} \emph{never}
processes an event that might later be preceded by an earlier incoming event.
Having an event $E_t$ at time $t$, the \ac{LP} must prove it will neve arrive
an event $E_{t'}$ with $t' < t$ on any of its communication channel.

\paragraph{Channel ordering and the importance of Lower Bounds:} Most
conservative protocols assume that messages on each channel are delivered in
non-decreasing timestamp order. Under this assumption, for each incoming channel
$c$, an \ac{LP} can maintain a lower bound $LB_c$ on the timestamp of the
next event that may arrive on $c$.
That said, if the \ac{LP} currently holds a candidate event with smallest
local timestamp $t_{\text{min}}$, then $t_{\text{min}}$ is \emph{safe} iff ~\cite{fujimoto1999Pads}:
\begin{equation}
	t_{\text{min}} \leq \min_{c \in \text{InputChannels}} \{LB_c\}
\end{equation}
The main challenge is therefore how to compute and propagate useful lower
bounds without introducing global barriers that limits parallelism.

\paragraph{Lookahead:} The ability of an \ac{LP} to predict a strictly
positive minimum time into the future before it can affect other \ac{LP}s
is called \emph{Lookahead}. If an \ac{LP} is at local simulation time
$t$ and can guarantee that any message it sends to a neighbour will have
timestamp at least $t + L$ for some $L > 0$, then it can advertise this
information as a lower bound. Lookahead often comes from model's structure, for
instance non-zero service times in queuing models or propagation latency in
network models. In general, the larger the lookahead, the easier and faster it
is for conservative protocols to find safe events and expose
parallelism~\cite{fujimoto1988LPDES}.

\paragraph{Blocking and Deadlock Avoidance:} in conservative execution, an
\ac{LP} may become blocked when its earliest local event is not provably safe.
In special cases where each \ac{LP} blocks for an earlier incoming message from
their input queues, a \emph{Deadlock} can occur.

In order to overcome such problems, the \ac{CMB} algorithm \newline
\cite{chandyMisra1983DDD, bryant1977simulation} avoids deadlock by exchanging
null messages. More specifically, each \ac{LP} periodically informs its
neighbours of a safe lower bound on future messages by sending to each outgoing
channel a null message advertising that it will not send any event with
timestamp smaller than $t + L$, where $t$ is current simulation time and $L$ is
the lookahead for that channel. Upon receiving null messages, an \ac{LP} can
update its lower bound on that channel $LB_c$.

Thanks to this synchronisation mechanism, deadlocks are avoided by paying an
extra communication overhead and potentially a large number of null messages
when lookahead is small.
\\
\newline

Overall conservative protocols are attractive because they never violate the
local causality constraint. However their performance depends critically on the
ability to compute good lower bounds (lookahead). When lookahead is poor,
\ac{LP}s spend a large fraction of time waiting, and communication overhead
could dominate execution.

\subsubsection{Speculative Execution}\label{sssec:speculative_pdes}

Speculative or Optimistic execution takes an opposite stance to the local causality
constraint with respect to conservative approaches, because rather than waiting until
an event is provably safe, an \ac{LP} processes events as soon as it can. This leads
to an increased level of concurrency, but causality violations could happen when
a \emph{straggler} event arrives. Optimistic protocols therefore include mechanisms
to \emph{detect} such violations and \emph{repair} the simulation state.

\paragraph{Time Warp and the Rollback procedure:}
The most influential optimistic approach is Jefferson's Time
Warp \cite{jefferson1987TWOS}, where each \ac{LP} maintains a local event list,
a local simulation clock (timestamp of most recently processed event $t_{LP}$), a
history of past states and an output log of messages it has sent. The last two
elements are used for enabling \ac{LP}'s \emph{rollback}.

The \ac{LP} repeatedly removes and processes its smallest timestamp event, even if it
is not known to be safe. If a new incoming event $E_{t'}$ with timestamp $t' < t_{LP}$,
$E_{t'}$ is marked as straggler events and the \ac{LP} performs a rollback as shown in
\Cref{fig:time_warp_rollbck}.
%
\includediagram{figures/diagrams/time_warp.tex}{Straggler event $E_{135}$
	arrival in an \ac{LP}, causing the \emph{rollback} at a consistent state
	with time $t < 135$, retrieved from the state log. Current simulation state
	is then $S_{130}$, and neighbors are advertised with respective
	anti-messages for events past $M_{130}$, namely $M_{140}$ and
	$M_{150}$.}{fig:time_warp_rollbck}
%
More specifically the \ac{LP}
%
\begin{enumerate*}[label=(\roman*)]
	\item restores the saved state corresponding or prior to time $t'$
	\item re-inserts any locally processed events with timestamp $t \geq t'$
	      back into the local event list and then
	\item "unsends" any messages previously sent with timestamp $t \geq t'$ by
	      sending \emph{anti-messages}, i.e. message cancellation notices.
\end{enumerate*}
%
Finally, the LP resumes forward execution reprocessing events in correct time
order.

\paragraph{Cascading Rollbacks, GVT and fossil collection:} by means of
anti-messages, Time Warp removes the effects of speculative execution from
other \ac{LP}s. If an \ac{LP} rolled back past a message it had sent, that
message must be undone at the receiver if it has already been processed. This
can trigger additional rollbacks, yielding a cascade. While we are assured that
no domino effect where the entire computation is eventually rolled back to the
beginning of the simulation~\cite{fujimoto1999Pads}, a practical implication is
that optimistic \ac{PDES} shifts some of the synchronisation cost from
\emph{waiting} (typical of conservative approaches) to work potentially done
twice.

Another important concern with speculative execution is that, in order to
support rollback, \ac{LP}s must retain state history and message logs, but
keeping unbounded history is infeasible. Therefore the concept of \ac{GVT} is
defined as a lower bound on the timestamp on any event that may still affect
the future of the simulation (shown in greyed queues' contents in
\Cref{fig:time_warp_rollbck}). Once \ac{GVT} has advanced to time $T$, no event
with timestamp smaller that $T$ can ever be rolled back again; hence, state
snapshots and logs older than $T$ can be safely discarded by the process of
\emph{fossil collection}.
\\
\newline

Summarising, optimistic execution tends to be effective where conservative
approaches fail, more specifically when dealing with dynamic dependencies
topology among \ac{LP}s with limited lookahead, when communication delays or
network latency are high (therefore waiting is too costly) and if causality
violations are rare.

\subsubsection{Conservative vs. Speculative}

While both approaches enforce the same correctness criterion, namely the
equivalence to a sequential timestamp-ordered \ac{DES}, the way they manage
uncertainty about future incoming events is fundamentally different:
\begin{itemize}
	\item Conservative approaches pays by \emph{stalling} and by exchanging
	      lower bounds (e.g. null messages) to avoid mistakes;
	\item Optimistic methods pays by \emph{repairing} mistakes through
	      rollbacks, anti-messages and periodic \ac{GVT}-based \emph{garbage
		      collection}.
\end{itemize}
As explained in \cite{Fujimoto1990PDES}, conservative approaches are faster in
presence of models which have a high degree of predictability, having good
lookahead and where the communication topology is static; conversely,
optimistic approaches are faster for dynamic topologies with irregular workload
distribution or models with poor lookahead. On the other hand conservative
approaches offer low memory overhead, whereas state-saving requirements for
optimistic approaches must be managed explicitly and could constitute a major
performance bottleneck.

Beyond the classical algorithmic taxonomy, a large body of work has focused on
making optimistic synchronisation practical on modern multi-core and
cluster-based platforms by reducing rollback costs and improving coordination,
proposing and evaluating \emph{runtime-level} mechanisms in optimistic
\ac{PDES}, ranging from engineering of rollback-oriented kernels and
APIs~\cite{pellegrini2012TROSCIPM} to OS and hardware-assisted techniques that
reduce wasted CPU work upon causality violations~\cite{silvestri2020IPI}.

Moreover, current research is going toward hybrid approaches (e.g. limited
optimism, window-based execution, adaptive throttling, \dots) as
\cite{piccione2023HSSPDES} which balances waiting and rollback depending on the
model structure and the target platform.

\subsection{Approximation-based Approaches}

Sometimes for really complex models, the application of standard approaches as
\ac{SSA} for complex reactions dynamics, could be infeasible despite what has
been seen for \ac{PDES}. In contexts where modelers require even faster
simulation speed, sometimes it can be obtained by making minor sacrifices in
simulation accuracy. In the following sections some notable approximation
techniques are presented, starting from the temporal approximation for large
stochastic reactions network, to methods that simplifies model's handling or
the model itself.

\subsubsection{Temporal Approximation}

The \ac{SSA} (\Cref{ssec:mdt}) is an exact method, but it is computationally
expensive because it treats every reaction as a discrete step. In systems with
large populations, the time between events becomes vanishingly small, halting
the simulation's progress on macroscopic scales \cite{Gillespie2001ACSSCRS}.

The $\tau$\emph{-Leaping} method accelerates execution by sacrificing exactness for
speed. It abandons the "exact time to next event" logic in favor of a
pre-selected time interval $\tau$ (the "leap"). The method relies on the
\emph{Leap Condition}, namely the assumption that $\tau$ is small enough that the
system's state, and consequently its propensity functions $a_j(x)$, remains
effectively constant during the interval $[t, t+\tau)$
\cite{Gillespie2001ACSSCRS}.

Under this assumption of local constancy, the number of times each reaction
$R_j$ fires during the interval is no longer dependent on the precise sequence
of other events. Instead, it can be approximated as a random sample from a
Poisson distribution:
%
\begin{equation}
	k_j \sim \mathcal{P}(a_j(x) \cdot \tau)
\end{equation}
%
The simulation cycle is thus reduced to:
\begin{enumerate}
	\item Select a time leap $\tau$ that satisfies the Leap Condition, i.e.
		ensuring propensities won't change "too much", for example by means of
		the $\tau$-selection procedure~\cite{yao2006ESSTLSM};
	\item For every reaction $R_j$, sample the number of firings $k_j$;
	\item Update the system state by applying the net effect of all $k_j$
		reactions simultaneously:
	      \begin{equation}
		      \mathbf{x}(t + \tau) = \mathbf{x}(t) + \sum_{j=1}^{M} k_j \boldsymbol{\nu}_j
	      \end{equation}
\end{enumerate}

This method also bridges the gap between discrete event simulation and
discrete time simulation (DTSS). When $\tau$ is small, it approaches the
accuracy of SSA; as $\tau$ grows, it approaches the behavior of a continuous
deterministic solver (like the Euler method) but with added stochastic noise
\cite{Gillespie2001ACSSCRS}.

\subsubsection{Scale Approximation}

A complementary approach addresses the dependency scale by dynamically
partitioning the system. In many multi-scale systems (e.g., biochemical
networks), some species are abundant while others are rare. Hybrid
solvers~\cite{haseltine2002ASCFSRSCK} treat the abundant species using
continuous differential equations (\ac{DES})—ignoring individual discrete
dependencies—while maintaining exact stochastic handling (\ac{DES}) for rare
species. This effectively reduces the size of the dependency graph by
"blurring" dense interaction webs into continuous rate functions.

\subsubsection{Structure Approximation}


\section{Motivation and Objectives}
In this chapter, a broad introduction of the simulation ecosystem is presented,
enabling the reader to understand the taxonomy of simulation systems and how
the problem of complex simulation systems is tackled, either by developing new
modeling and simulation formalisms such as agent-based simulation
or by designing clever algorithms by studying complex scenarios such as
biochemical systems. As stated at the beginning of this section, every
approach presented involves some kind of dependency management

